% Reply to Reviewer 1, Q3 (noise-free gap; restrict "on par with an oracle").
% Numbers: figures/flyvis/cv_table_known_ode_vs_gnn_fullsample.tex (Tab. 1 twin).
% V_rest row is the FULL-SAMPLE value (see Q2); cluster and R2_W are unfiltered.

\paragraph{Q3.}
We restrict the claim. ``On par with an oracle'' will be stated for
$\sigma > 0$, and the noise-free row will be presented as what it is, the regime
where the inverse problem is ill-posed.

\begin{center}
{\footnotesize
\begin{tabular}{lcccc}
\toprule
 & GNN $\sigma = 0$ & Known ODE $\sigma = 0$ & GNN $\sigma = 0.05$ & Known ODE $\sigma = 0.05$ \\
\midrule
$R^2_{\widehat{W}}$                                & $0.89$  & $0.96$ & $0.99$ & $0.99$ \\
$R^2_{\widehat{V}^{\mathrm{rest}}}$ (full-sample)  & $-0.39$ & $0.86$ & $0.76$ & $0.95$ \\
cluster acc.                                       & $0.86$  & $0.92$ & $0.92$ & $0.93$ \\
\bottomrule
\end{tabular}}
\end{center}

\textbf{Why the gap is there.} Noise-free, Appx.~C gives $115{,}223$ near-null
directions ($26.5\%$ of edges). Along those directions the estimate is set by the
prior rather than by the data, and the oracle's prior is the true equation with
$\ell_1/\ell_2$ on its own parameters, while the GNN must identify $f_\theta$ and
$g_\phi$ as well. It therefore pays more for the same degeneracy. Process noise
removes the degeneracy and with it the gap.

\textbf{On the premise.} $\sigma$ is process noise. We have not calibrated its
level against biology and do not claim to know it, but transmission from neuron
to neuron is stochastic: release is probabilistic, channels are noisy, and
membrane potentials fluctuate under background drive. So the open question is
where real circuits sit on the $\sigma$ axis, not whether $\sigma = 0$; the
noise-free setting is the artificial one, a knob we can turn off in silico and
nothing turns off in vivo. We report both ends of the axis, $\sigma = 0$ and
$\sigma = 0.5$, so that a reader can bracket any real value. What real recordings
add instead is
measurement noise $\gamma$, and there we agree the method is not yet usable
(W4): $R^2_{\widehat{W}} = 0.63$ at $\gamma = 0.1$ for the GNN and $0.69$ for the
oracle. The binding constraint for real data is $\gamma$, not $\sigma = 0$, which
is why we do not treat the noise-free gap as the headline.

We keep the noise-free row as evidence for the paper's thesis rather than as a
result to defend. One-step $r = 1.00$ with $R^2_{\widehat{W}} = 0.89$ and
$R^2_{\widehat{V}^{\mathrm{rest}}} = -0.39$ over all neurons is prediction
without mechanism, which is what the identifiability analysis predicts and what
the paper is about.
